\section{本讲小结}
磁通量的正负来自曲面法线，电动势的正负来自与它配套的回路方向。先选方向再求导，比在计算末尾补方向可靠。对不均匀场要积分有向分量，不能只比较面积。

法拉第定律确定闭合回路的总电动势；局部运动导体用 $\int(\bm u\times\bm B)\cdot\dd\bm l$，轴对称感生场由环流求得。电动势、电流和端电势差是不同的物理量。粒子在感生场中的运动要分别列切向与径向方程，电场功负责动能变化，磁力与约束共同决定弯曲轨迹。

多个因素同时变化时，对有向磁通完整求导，检查动生与感生项是否相消。电量、热量与电流峰值分别对应一次积分、平方积分及极值；它们不能互相代替。弯曲导线作刚性平移时，电动势取决于端点投影，电阻取决于实际长度，这两个几何量必须分开。

